\documentclass[7pt,landscape]{article}
\usepackage[utf8]{inputenc}
\usepackage[margin=0.18in]{geometry}
\usepackage{multicol}
\usepackage{titlesec}
\usepackage{enumitem}
\usepackage{xcolor}

\definecolor{navy}{RGB}{0, 51, 102}
\definecolor{linuxc}{RGB}{160, 88, 24}
\definecolor{winc}{RGB}{45, 95, 140}
\setlength{\columnsep}{0.25in}
\setlist{nosep, leftmargin=10pt}
\titlespacing*{\section}{0pt}{1pt}{0.5pt}
\setlength{\parindent}{0pt}
\linespread{0.90}

\titleformat{\section}{\small\bfseries\color{navy}}{\thesection}{0.5em}{}[\titlerule]
\pagestyle{empty}

\begin{document}
\textbf{Role Reference: Captain} \hfill \textit{Last Updated: Aug 2026}
\raggedcolumns
\begin{multicols*}{3}

\section{ROLE RESPONSIBILITIES}
Injects, phone, flex support, coordinating the team. Rarely fixes anything directly. Makes sure every fix has an owner and nothing falls through.

\section{ORDER OF OPERATIONS}
\textbf{Round start:} Confirm every role staffed. Confirm shared log/template open. \\
\textbf{Inject arrives:} Read fully. Note deadline. Assign owner immediately. \\
\textbf{Foothold reported:} Get host, time, symptom. Assign loop: threat hunter continues user/hole, SIEM lead checks other hosts, engineer stages fix. \\
\textbf{Findings converge:} Confirm engineer has what they need. Confirm scope logged before eviction. \\
\textbf{Verified fixed:} Log closed, timestamped, with confirmer. Move to next item.

\section{THINK LIKE RED TEAM}
\begin{itemize}
    \item Weigh a fix by how much of red team's playbook it disables, not just whether it closes today's hole.
    \item Push the engineer toward removals that cost red team more than blue team or the scored service: unused interpreters, unused remote-management protocols, reused local admin passwords.
    \item A quiet stretch can mean red team is regrouping around tooling you removed, not that they've given up.
\end{itemize}

\section{TRIAGE PRIORITY}
\begin{itemize}
    \item Severity first: scoring/safety impact now, or can it wait.
    \item Effort second: 5-minute fix ahead of a 30-minute one, all else equal.
    \item An owner, always. Untracked work does not happen.
    \item A visible queue. Everyone sees what is pending, not just you.
\end{itemize}

\section{THE TRACKING TEMPLATE}
Keep running the entire round, not reconstructed after.
\begin{itemize}
    \item Time detected, reported by, host(s) affected
    \item User/account involved, entry vector (the hole)
    \item Lateral movement: yes/no, where
    \item Exfiltration: yes/no, what
    \item Owner assigned, actions taken (timestamped)
    \item Verified closed: time, by whom
\end{itemize}

\section{HANDLING INJECTS}
\begin{itemize}
    \item Read the whole queue before assigning from it.
    \item Schedule injects alongside technical work, not after it.
    \item Do not pull a threat hunter or engineer off active eviction for a routine inject. Use flex, or handle it yourself.
    \item Keep a standard inject response template ready: restate the ask, list what was done, note timestamps. Saves time under a deadline.
\end{itemize}

\section{RUNNING THE COMMS CHANNEL}
\begin{itemize}
    \item One channel for status updates, separate from casual chat, if your tool allows it.
    \item Pin the tracking template link and the current priority list.
    \item Use short, structured updates: role, host, status, next step. Not paragraphs.
    \item A 60-second radio check every 15-20 minutes catches silent blockers early.
\end{itemize}

\section{YOUR OWN CHECKS}
Things you can verify yourself without pulling a specialist off task.
\begin{itemize}
    \item \texttt{ping [host]} / \texttt{Test-Connection}: is it even up.
    \item \texttt{curl -I [url]} or a browser: is the scored service responding.
    \item \texttt{ssh [host]} or RDP: can you still get in at all.
    \item \texttt{whoami}: confirm you did not just get logged in as someone unexpected.
\end{itemize}

\section{DELEGATION MAP}
\begin{itemize}
    \item \textit{Domain Admin:} Anything AD-specific.
    \item \textit{Threat hunters:} Detect, find user, find hole, assess scope, per OS.
    \item \textit{Engineer(s):} Evict and patch, after handoff.
    \item \textit{SIEM lead:} Scope across every host. Watches for return.
\end{itemize}

\section{FLEX BRIEFING: DOMAIN ADMIN}
Enough to cover if nobody else can.
\begin{itemize}
    \item Baseline: \texttt{Get-ADUser -Filter *}, \newline \texttt{Get-ADGroupMember "Domain Admins"}.
    \item Health: \texttt{dcdiag}, \texttt{repadmin /replsummary}.
    \item Policy: \texttt{Get-GPO -All}, \texttt{gpresult /r} to confirm it applied.
    \item Red flag: a new Domain Admin, or unusual Kerberos ticket activity.
    \item Reset: \texttt{Set-ADAccountPassword}, \newline or \texttt{Disable-ADAccount} if unsure.
\end{itemize}

\section{FLEX BRIEFING: THREAT HUNTER}
\begin{itemize}
    \item Who's logged in: \texttt{w}/\texttt{who} \textcolor{linuxc}{(Linux)}, \texttt{query user} \textcolor{winc}{(Windows)}.
    \item Who logged in recently: \texttt{last} \textcolor{linuxc}{(Linux)}, Event ID 4624/4625 \textcolor{winc}{(Windows)}.
    \item Trace a process: \texttt{ps -o user,pid,ppid,cmd -p [PID]} \textcolor{linuxc}{(Linux)}, \texttt{Get-CimInstance Win32\_Process} for parent PID \textcolor{winc}{(Windows)}.
    \item Find the hole: check auth logs for the login window around the suspicious timestamp.
    \item Scope: \texttt{ss -tunp} / \texttt{Get-NetTCPConnection} for unfamiliar connections.
\end{itemize}

\section{FLEX BRIEFING: ENGINEER}
\begin{itemize}
    \item Kill it: \texttt{kill -9 [PID]} \textcolor{linuxc}{(Linux)}, \texttt{Stop-Process -Id [PID] -Force} \textcolor{winc}{(Windows)}.
    \item Remove persistence: cron/\texttt{/etc/cron.d} \textcolor{linuxc}{(Linux)}, Task Scheduler/Run keys \textcolor{winc}{(Windows)}.
    \item Patch the actual vulnerability, not just the symptom. Rotate anything possibly exposed.
    \item Protect a fix: \texttt{chattr +i} on \textcolor{linuxc}{Linux}, tighter \texttt{icacls} on \textcolor{winc}{Windows}.
    \item Confirm the scored service is still up after any change.
    \item If time allows: strip unused interpreters (Python, Perl), disable unused WinRM/PSRemoting.
\end{itemize}

\section{FLEX BRIEFING: SIEM LEAD}
\begin{itemize}
    \item Health: \texttt{so-status}. Access: \texttt{so-allow} if the web UI refuses you.
    \item Use Hunt for ad-hoc queries, not just the prebuilt Alerts view.
    \item Lateral movement: \texttt{event.dataset:conn} filtered to the host in question.
    \item Other hosts: search Suricata alerts by the same signature across every sensor.
    \item If Wazuh is deployed, File Integrity Monitoring alerts confirm a specific file changed fast.
\end{itemize}

\section{SCOREBOARD}
\begin{itemize}
    \item Points typically split between service uptime (checked on an interval) and injects (graded on submission quality and timeliness).
    \item Uptime is scored automatically and continuously; a service down for ten minutes costs more than one incident report submitted five minutes late.
    \item Know which services are actually scored vs. which are just present. Do not spend a scarce engineer on a box that earns nothing.
    \item Ask White Team early if scoring criteria for an inject are ambiguous. A clarifying question costs less than a wrong guess.
\end{itemize}

\section{SIMULTANEOUS INCIDENTS}
\begin{itemize}
    \item Rank by scored-service impact first, safety second, everything else after.
    \item Do not let a loud, obvious problem eclipse a quiet one on a more valuable host. Check the full host list, not just what people are shouting about.
    \item Assign a distinct owner per incident even if hands are shared across them; the tracking template still needs one line per incident, not one merged mess.
    \item If truly overloaded, tell the team explicitly what is being deprioritized and why, out loud, not silently.
\end{itemize}

\section{WORKING WITH WHITE TEAM / JUDGES}
\begin{itemize}
    \item Report an outage or scoring dispute the moment it is noticed, not after trying to fix it quietly first.
    \item Keep language factual: what you observed, when, what you believe caused it. Save the argument for after the round.
    \item A formal timeout or scoring pause exists for a reason. Use it for a genuine platform issue, not as a stalling tactic.
    \item Written communication (chat, ticket) beats verbal when a dispute might need to be revisited later.
\end{itemize}

\section{DEBRIEF}
\begin{itemize}
    \item Walk the timeline from the template, not memory.
    \item Find what slowed the loop, not who to blame.
    \item Ask where a handoff was slow or unclear. That is usually the real fix.
\end{itemize}

\end{multicols*}
\end{document}